\documentclass[UTF8,fontset=fandol,a4paper,11pt]{ctexart}
\usepackage[margin=22mm,top=23mm,bottom=22mm,headheight=15pt]{geometry}
\usepackage{amsmath,amssymb,bm,mathtools,booktabs,longtable,tabularx,array,multirow}
\usepackage[table]{xcolor}
\usepackage{graphicx,tikz}
\usetikzlibrary{arrows.meta,positioning,fit,calc}
\usepackage[most]{tcolorbox}
\usepackage{enumitem,fancyhdr,lastpage,listings,xurl,placeins,needspace,float}
\usepackage[numbers,sort&compress]{natbib}
\usepackage[unicode,colorlinks=true,linkcolor=navy,urlcolor=tealblue,citecolor=tealblue]{hyperref}
\definecolor{navy}{HTML}{18344D}
\definecolor{tealblue}{HTML}{087F8C}
\definecolor{pale}{HTML}{EDF5F7}
\definecolor{warm}{HTML}{9A5A18}
\definecolor{ink}{HTML}{253443}
\hypersetup{pdftitle={SuperMimic 论文与代码对齐报告},pdfauthor={代码与文献审读},pdfsubject={arXiv:2610.28480v1 与 supermimic-vibed 8bfc1158 的实现对齐}}
\ctexset{section={format=\Large\bfseries\color{navy}},subsection={format=\large\bfseries\color{navy}},subsubsection={format=\normalsize\bfseries}}
\setlength{\emergencystretch}{3em}
\setlength{\parskip}{0.35em}
\linespread{1.16}
\setcounter{tocdepth}{2}
\setcounter{secnumdepth}{3}
\setlist[itemize]{leftmargin=1.6em,itemsep=0.15em,topsep=0.3em}
\setlist[enumerate]{leftmargin=2em,itemsep=0.2em,topsep=0.3em}
\renewcommand{\arraystretch}{1.22}
\setlength{\tabcolsep}{5pt}
\newcolumntype{Y}{>{\raggedright\arraybackslash}X}
\newcolumntype{P}[1]{>{\raggedright\arraybackslash}p{#1}}
\pagestyle{fancy}
\fancyhf{}
\fancyhead[L]{\small\color{navy}SuperMimic / 论文与代码对齐}
\fancyhead[R]{\small\color{navy}源码快照 8bfc1158}
\fancyfoot[L]{\footnotesize\color{gray}非官方复现审读 · 2026-10-30}
\fancyfoot[R]{\small\thepage}
\renewcommand{\headrulewidth}{0.35pt}
\newcommand{\code}[1]{\nolinkurl{#1}}
\newcommand{\src}[3]{\href{https://github.com/zita-ch/supermimic_vibed/blob/8bfc1158bd968f729fcc65d5e897a05b05ea3d36/#1\#L#2-L#3}{\nolinkurl{#1}:#2--#3}}
\newcommand{\paper}[1]{\textbf{论文定位：}#1\cite{liao_supermimic_2026}。}
\newcommand{\source}[1]{\par\noindent{\small\color{tealblue}\textbf{源码定位：}#1。}\par}
\newcommand{\R}{\mathbb{R}}
\newcommand{\E}{\mathbb{E}}
\newcommand{\sg}{\operatorname{sg}}
\newcommand{\dd}{\mathrm{d}}
\newcommand{\heading}{\operatorname{yaw}}
\newcommand{\six}{\operatorname{six}}
\newtcolorbox{insight}[1]{enhanced,colback=pale,colframe=tealblue,boxrule=0.55pt,arc=1.5mm,title={#1},fonttitle=\bfseries,coltitle=white,before skip=8pt,after skip=8pt}
\newtcolorbox{caution}[1]{enhanced,colback=orange!4,colframe=warm,boxrule=0.55pt,arc=1.5mm,title={#1},fonttitle=\bfseries,coltitle=white,before skip=8pt,after skip=8pt}
\lstset{language=Python,basicstyle=\ttfamily\footnotesize,keywordstyle=\color{navy}\bfseries,commentstyle=\color{gray},stringstyle=\color{tealblue},breaklines=true,columns=fullflexible,keepspaces=true,showstringspaces=false,frame=single,rulecolor=\color{gray!30},backgroundcolor=\color{gray!3},aboveskip=0.6em,belowskip=0.6em}
\begin{document}
\hypersetup{pageanchor=false}
\begin{titlepage}
\thispagestyle{empty}
\begin{tikzpicture}[remember picture,overlay]
\fill[navy] (current page.north west) rectangle ([yshift=-82mm]current page.north east);
\fill[tealblue] ([yshift=-82mm]current page.north west) rectangle ([yshift=-84mm]current page.north east);
\end{tikzpicture}
\vspace*{2mm}
{\color{white}\sffamily\large PAPER $\leftrightarrow$ CODE\par}
\vspace{7mm}
{\color{white}\Huge\bfseries SuperMimic\par}
\vspace{3mm}
{\color{white}\LARGE 论文与代码对齐报告\par}
\vspace{28mm}
{\Large\bfseries 从可微世界模型到全身运动跟踪\par}
\vspace{4mm}
{\large 框架、张量接口、核心公式、梯度路径与复现边界\par}
\vspace{9mm}
\noindent\begin{tabularx}{\textwidth}{@{}P{27mm}Y@{}}
\toprule
论文 & \textit{SuperMimic: Scaling Humanoid Whole-Body Control with One Differentiable World Model}\\
作者 & Qiayuan Liao, Yiyang Shao, Junfeng Long, Kevin Zakka, Roberto Horowitz, Pieter Abbeel, Koushil Sreenath\\
文献版本 & arXiv:2610.28480v1，2026-10-20；Zotero 全文，共19页\\
代码对象 & \url{https://github.com/zita-ch/supermimic_vibed}\\
固定提交 & \texttt{8bfc1158bd968f729fcc65d5e897a05b05ea3d36}\\
报告日期 & 2026年10月30日（Asia/Shanghai）\\
\bottomrule
\end{tabularx}
\vspace{8mm}
\begin{insight}{报告结论先读}
该仓库实现了 SuperMimic 的核心三阶段训练循环：仿真采样、世界模型拟合、基于想象轨迹的 BPTT 策略与价值学习。它是面向 LAFAN1 的\textbf{非官方、部分范围复现}，不能直接视为论文全部实验的官方代码。本文将论文规定、代码事实、解释性推导和已运行验证明确区分。
\end{insight}
\vfill
{\small 交付内容：可检索 PDF、可编辑 LaTeX 工程、固定版本来源清单及 CPU 验证脚本。\par}
{\small 本次未进行 GPU 训练或实机实验；性能数据均注明出处与证据等级。\par}
\end{titlepage}
\hypersetup{pageanchor=true}
\pagenumbering{roman}
\tableofcontents
\clearpage
\pagenumbering{arabic}
\section{阅读结论与证据边界}
\subsection{这份代码究竟复现了什么}
SuperMimic 的关键并不是“换一个更大的模仿网络”，而是改变策略梯度的获得方式：用普通刚体仿真器收集状态转移，学习可微的动力学替代模型，再把跟踪回报沿短想象轨迹反向传播到策略。仿真器负责生成真实训练转移，\textbf{不需要提供梯度}；策略学习所需的雅可比来自学习到的世界模型\cite{liao_supermimic_2026}。

目标仓库自述为独立、非官方复现，并明确指向论文图4(a)的受控 LAFAN1 比较\cite{zhang_supermimic_vibed}。代码默认载入40段 G1 重定向 LAFAN1 动作，README 将总时长记为2.45小时；这不等于论文主实验所用264.56小时、129,514段筛选后 BONES-SEED 动作。论文的多接口控制、扩展规律、双 GPU 训练与实机展示也没有在该仓库形成相应完整实现。

\begin{insight}{最重要的三项判断}
\begin{enumerate}
\item \textbf{方法主干对齐。} 代码确有在线数据收集、$k=8$步自回归动力学学习、$H=32$步可微想象、软存活权重、终端价值 bootstrap 与固定的 $\lambda$-return 目标。
\item \textbf{实现不等于逐字复刻。} 代码选择了局部状态增量、在线标准化、额外6维相对朝向观测、解析辅助头、sigmoid 软终止、bf16 矩阵乘法和额外预热等。需要说明各自来源，而不能把全部细节倒推为论文原设定。
\item \textbf{跑通局部逻辑不等于复现性能。} 本次完成静态审读及16项 CPU 检查；没有执行 Isaac Lab 仿真、GPU 训练或检查已训练 checkpoint。仓库的0.829回报及40/40成功是作者记录，不是本次测量结果。
\end{enumerate}
\end{insight}

\subsection{版本、文献与四种证据标签}
论文来自用户 Zotero 条目 \code{IQ8SWVH7}，读取的 PDF 附件为 \code{TQR99565}，正文首页标注 \code{arXiv:2610.28480v1}，共19页。BibTeX 引用键 \code{liao_supermimic_2026} 与上述 Zotero 条目键不是同一种标识。PDF 的 SHA-256 和代码文件校验和保存在随附 \code{source_manifest.json} 中。

源码固定在封面列出的完整提交；正文链接均指向该提交的具体行，而不是可变的 \code{main} 分支。仓库中“当时官方代码未发布”的叙述只作为其历史背景，\textbf{本报告未核验报告日期时官方仓库是否已更新}。

\noindent\begin{tabularx}{\textwidth}{@{}P{20mm}Y Y@{}}
\toprule
证据类别 & 本文使用方式 & 不应推出的结论\\
\midrule
论文规定 & 章节、公式、附录表和页码定位 & 不能证明非官方实现完全遵守\\
代码事实 & 固定提交中的函数、数据流和参数 & 不能仅由源码证明训练成绩\\
解释推导 & 由公式及代码推得的张量维数、梯度路径和影响 & 不替代消融实验与收敛证明\\
运行验证 & 明确列出执行环境、检查项和通过数 & 不外推到物理仿真、CUDA 或实机\\
\bottomrule
\end{tabularx}

\subsection{建议的阅读顺序}
建议先读第\ref{sec:framework}节的数据流及第\ref{sec:wm}、\ref{sec:actor}节的核心公式，再核对第\ref{sec:interface}节接口。第\ref{sec:matrix}节集中列出对齐差异，第\ref{sec:evaluation}节解释实验可比性。

\FloatBarrier
\section{系统框架：两个前向世界，一条策略梯度链}\label{sec:framework}
\paper{第3.1节、图2、附录D算法A.1（第4、15页）}
\subsection{三阶段训练循环}
\begin{figure}[htbp]
\centering
\begin{tikzpicture}[font=\small,node distance=8mm and 8mm,box/.style={draw=tealblue,rounded corners,fill=pale,align=center,text width=43mm,minimum height=17mm},arr/.style={-{Latex[length=2mm]},thick,draw=navy}]
\node[box] (sim) {刚体仿真采样\\$a=\pi_\theta(o,g)+\epsilon$};
\node[box,right=of sim] (buf) {GPU replay buffer\\$(s,a,s',\xi,f,\ldots)$};
\node[box,right=of buf] (fit) {自回归动力学拟合\\更新 $\omega$，$k=8$};
\node[box,below=14mm of buf] (imag) {冻结世界模型参数\\$H=32$步想象与解析头};
\node[box,left=of imag] (actor) {最大化想象回报\\BPTT 更新 $\theta$};
\node[box,right=of imag] (critic) {固定 $\lambda$-return 目标\\监督更新 $\psi$};
\draw[arr] (sim)--(buf);
\draw[arr] (buf)--(fit);
\draw[arr] (buf)--node[right,font=\scriptsize]{抽起点及历史}(imag);
\draw[arr] (fit)--(imag);
\draw[arr] (imag)--(actor);
\draw[arr] (imag)--(critic);
\draw[arr] (actor)--(sim);
\draw[arr,dashed] (critic.north west)--node[above,sloped,font=\scriptsize]{终端价值}(imag.north east);
\end{tikzpicture}
\caption{根据论文图2与固定版本代码重新绘制的训练数据流。虚线表示价值 bootstrap 的依赖，不表示对仿真器求导。}
\end{figure}

\textbf{阶段一：Collect。} 当前策略依据带噪观测与未来动作目标输出29维动作。采样时加入标准差0.1的高斯噪声，执行后的动作写入缓存。仿真器产生下一状态，解析 MDP 代码为其计算回报和终止标签。旧策略产生的数据仍可用于动力学监督，因为物理转移不因“谁采样了它”而改变。

\textbf{阶段二：Fit。} 从缓存取2048个8步窗口。世界模型只在窗口第一帧使用真实状态，之后以自身预测为输入，依次执行缓存中的实际动作。每次外层迭代做10次世界模型优化，而非只做一次单步回归。

\textbf{阶段三：Improve。} 从缓存取4096个合法起点及其10帧观测历史，使用当前策略和固定的世界模型进行32步想象。解析计算回报与软终止，累加折扣存活回报并加末端 critic 值；对策略做一次反向传播。随后固定本批想象轨迹产生的价值目标，对 critic 做5轮、每轮4个 minibatch 的拟合。
\source{\src{supermimic/train.py}{266}{321}；\src{supermimic/train.py}{147}{156}；\src{supermimic/train.py}{175}{252}}

\subsection{模块分工与调用关系}
\noindent\begin{tabularx}{\textwidth}{@{}P{32mm}Y Y@{}}
\toprule
模块 & 主要责任 & 核心调用或输出\\
\midrule
\code{train.py} & 调度采样、模型、actor、critic，管理优化器与检查点 & \code{collect / wm_update / policy_update}\\
\code{env.py} & Isaac Lab/物理后端封装，动作执行、重置与逐环境上下文 & \code{TrackingSim.step/reset/state}\\
\code{motion.py} & CSV 重采样、参考速度、跟踪体缓存与未来目标 & \code{MotionLib.goal/ref_state}\\
\code{mdp.py} & 状态编解码、观测、回报、硬/软终止与 critic 特征 & 仿真与想象共享同一套解析语义\\
\code{networks.py} & 世界模型、策略编码器/解码器、critic、在线归一化 & \code{WorldModel / Policy / Critic}\\
\code{buffer.py} & GPU 环形缓存、跨终止掩码、历史恢复 & 时间 $\times$ 环境的稠密存储\\
\code{kinematics.py} & 可微 FK 与 link-origin 速度递推 & 14个跟踪体的位姿和速度\\
\code{robot.py} & G1关节顺序、PD参数、动作缩放与资产路径 & 29自由度与14个跟踪体定义\\
\code{evaluate.py} & 从参考首帧确定性运行并汇总指标 & clip-mean / frame-weighted 回报与失败数\\
\bottomrule
\end{tabularx}

\subsection{训练与部署的信息边界}
世界模型接收完整状态和仿真上下文 $\xi$，critic 还可接收参考误差等特权量；策略则接收带噪观测历史与参考目标，不直接接收摩擦、质心偏移等真实参数。这种不对称输入让训练侧利用更多信息，而控制器仍依据可构造的观测动作。

世界模型既不读取未来参考目标，也不把 motion ID 当作物理参数；动作参考只决定策略行为、回报和终止，因此一个动力学模型可以共享于很多动作。评估时只需要策略和解析观测/目标构造，\textbf{不运行世界模型和 critic}。这份仓库没有真实机器人通信和部署执行器，不能把“评估只加载策略”理解为已经实现了实机部署。

\FloatBarrier
\section{MDP 与张量契约：先对齐量，再对齐网络}\label{sec:interface}
\paper{附录A、表A.1--A.3（第13--14页）及附录C表A.4（第15页）}
\subsection{状态、观测、目标与上下文不是同一个向量}
代码的内部状态为
\begin{equation}
 S=(p,R,v,\omega,q,\dot q),\quad p,v,\omega\in\R^3,\ R\in SO(3),\ q,\dot q\in\R^{29}.
\end{equation}
$p,R$ 是 pelvis 的世界位姿，$v,\omega$ 是该 link 原点的世界线速度和角速度，不是把质心速度无条件当作 link 速度。缓存把旋转矩阵转为 \code{wxyz} 四元数；外部 CSV 与 Isaac Lab 状态的四元数则按 \code{xyzw} 转换。把这两种约定混用会直接破坏参考动作和模型监督。

\begin{table}[htbp]
\centering\small
\noindent\begin{tabularx}{\textwidth}{@{}P{36mm}P{27mm}Y@{}}
\toprule
对象 & 默认形状（批大小$B$） & 含义与计算\\
\midrule
缓存状态 & $B\times71$ & $3+4+3+3+29+29$，用四元数存旋转\\
世界模型不变特征 & $B\times71$ & $1+6+3+3+29+29$，不是缓存状态的别名\\
世界模型输入 & $B\times133$ & 不变特征71 + 动作29 + 上下文33\\
动力学增量输出 & $B\times70$ & $3+3+3+3+29+29$，旋转在三维切空间预测\\
单帧策略观测 & $B\times99$ & 论文93维 + 额外6维相对根朝向\\
观测历史 & $B\times10\times99$ & 展平为990维，最早帧在前\\
单帧参考特征 & $B\times73$ & $3+6+3+3+29+29$\\
未来目标 $g$ & $B\times730$ & $K=10$帧参考特征展平\\
隐变量 $z$ / 动作 & $B\times32$ / $B\times29$ & $z$ 经 $\tanh$，动作输出无 $\tanh$\\
策略解码器输入 & $B\times1022$ & 990维历史 + 32维隐变量\\
critic 输入 & $B\times914$ & 不变特征71 + 额外113 + 目标730\\
\bottomrule
\end{tabularx}
\caption{由固定版本代码计算的默认接口维数。符号“状态维数71”在两处出现，但量的语义不同。}
\end{table}
\source{\src{supermimic/dimensions.py}{1}{2}；\src{supermimic/mdp.py}{54}{81}；\src{supermimic/train.py}{69}{78}}

\subsection{观测与动作：哪些属于论文，哪些是额外选择}
论文表A.1给出93维本体观测：$q-q^0$（29）、$\dot q$（29）、机体系角速度（3）、投影重力（3）、上一动作（29），并使用10帧历史。代码实际构造
\begin{equation}
\begin{split}
 \bar o_t=\big[&q_t-q^0-\Delta q^0,\ \dot q_t,\ R_t^\top\omega_t,\ -R_t^\top e_3,\ a_{t-1},\\
 &\six(R_t^\top R_t^{\rm ref})\big]\in\R^{99}.
\end{split}
\end{equation}
最后6维是\textbf{额外的参考相对根朝向}，原论文表A.1未列出。这里 $\six$ 是矩阵前两列按代码顺序展平，不是欧拉角。它显式输入跟踪朝向误差；不能因为不包含全局位置，就称为与论文观测完全相同。

代码按表A.1对前四类量加入独立均匀噪声，半宽依次为0.01 rad、0.5 rad/s、0.2 rad/s、0.05；上一动作不加噪。新增6维也加入半宽0.05的分量噪声。之后 \code{Policy.obs_norm} 再做在线均值方差标准化。\code{mdp.py} 虽定义 \code{OBS_SCALE}，但当前 \code{obs()} 和策略前向并未使用该固定缩放常量，不能把它误写成实际观测预处理。

动作是归一化关节设定量，不是网络直接输出电机力矩：
\begin{equation}
 q^{\rm sp}=q^0+\Delta q^0+\alpha\odot a,\quad
 K_p=I_a\omega_n^2,\quad K_d=2\zeta I_a\omega_n,\quad
 \alpha_j=\frac{0.25\,\tau^{\max}_j}{K_{p,j}},
\end{equation}
其中 $I_a$ 为反射电枢惯量，$\omega_n=2\pi\cdot10$，$\zeta=2$。机器人配置使用隐式PD驱动，并设置 effort/velocity 限制。策略解码器是线性输出；采样时对\emph{加噪后的}动作裁剪到 $[-20,20]$，这不是通常意义的 $[-1,1]$ 动作空间，想象与评估中的策略输出也没有这一步裁剪。
\source{\src{supermimic/mdp.py}{18}{35}；\src{supermimic/mdp.py}{94}{105}；\src{supermimic/robot.py}{26}{36}；\src{supermimic/robot.py}{88}{91}；\src{supermimic/env.py}{192}{205}；\src{supermimic/train.py}{274}{280}}

\subsection{参考动作管线与未来目标的具体构造}
CSV 每帧包含根位置3、根四元数4、29个关节角，共36个数，关节顺序严格按照 \code{G1_JOINT_NAMES}。\code{MotionLib} 从30 Hz重采样到50 Hz：根位置与关节角线性插值，四元数做最短弧 SLERP；线速度与关节速度用差分，角速度用相邻旋转的 $SO(3)$ 对数计算。随后通过 FK 预计算14个跟踪体的位姿和速度。

令 $H_f=\heading(R_f^{\rm ref})$，未来帧 $j=f+1,\ldots,f+K$ 的特征为
\begin{equation}
 g_{f,j}=\left[
 H_f^\top(p_j^{\rm ref}-p_f^{\rm ref}),\ {}
 \six(H_f^\top R_j^{\rm ref}),\ {}
 (R_j^{\rm ref})^\top v_j^{\rm ref},\ {}
 (R_j^{\rm ref})^\top\omega_j^{\rm ref},\ {}
 q_j^{\rm ref},\ 0.1\dot q_j^{\rm ref}
 \right].
\end{equation}
目标的相对位移/朝向以\textbf{当前参考 heading} 为基准，速度则以\textbf{各自未来参考根坐标系} 表示；不要把它们统称为同一机体系。$K=10$提供前方0.02--0.20秒的参考，编码器接收730维扁平向量，而非 Transformer 序列。

查询超出 clip 尾部时，位置、朝向与关节角保持最后一帧，合成延长部分的速度置零；真实最后一帧原本的速度仍保留。\code{goal()} 不修改源参考数据，也不修改用于回报的 \code{ref_state()}。这一区别避免“为了补齐网络输入，把真实尾帧也静止化”。
\source{\src{supermimic/motion.py}{24}{70}；\src{supermimic/motion.py}{94}{120}；\src{supermimic/robot.py}{39}{58}}

\subsection{回报：同一解析函数同时服务仿真和想象}
设 $\mathcal T$ 为14个跟踪体。代码先计算跟踪误差的六个标量：各身体的平均位置平方误差 $E_p$、平均旋转测地角平方 $E_R$、平均线/角速度平方误差 $E_v,E_\omega$，以及根位置和根朝向误差 $E_{rp},E_{rR}$。
\begin{equation}\label{eq:reward}
\begin{split}
 r_{t+1}=\Delta t\Big[&e^{-E_p/0.3^2}+e^{-E_R/0.4^2}
 +e^{-E_v/1.0^2}+e^{-E_\omega/3.14^2}\\
 &+0.5e^{-E_{rp}/0.3^2}+0.5e^{-E_{rR}/0.4^2}
 -0.1\lVert a_t-a_{t-1}\rVert^2-2P_{\rm lim}(q_{t+1})\Big],
\end{split}
\end{equation}
其中 $\Delta t=0.02$ 秒；$P_{\rm lim}$ 是越过关节上下限的线性 ReLU 和，不是平方惩罚。代码用关节范围中间90\%作为软限制。六个正奖励的权重和为5，因此 \code{MAX_STEP_REWARD}=0.1，供日志和评估归一化使用。

\textbf{坐标重锚定需要逐项读。} 身体位置/朝向的参考先按机器人当前根的xy位置与 yaw 重锚定，锚点高度仍使用参考根高度；根位置和根朝向两项仍比较原始参考全局量。身体速度两项使用原始世界坐标参考速度，没有同步做上述 yaw 重锚定。因此它不是“所有奖励均对全局平移与旋转不变”的实现；这也是检查奖励重构时不能省略的细节。原论文表A.2列出了误差形式，但未把这一重锚定实现细节完整展开，本报告只将其记为代码事实。
\source{\src{supermimic/mdp.py}{84}{89}；\src{supermimic/mdp.py}{108}{145}}

\subsection{终止、随机化与解析 FK}
仿真硬终止是
\begin{equation}\label{eq:hardterm}
 d_{t+1}=\mathbf1\{\lvert p_z^{\rm ref}-p_z\rvert>0.25\ \lor\ {}
 \lvert(R^{\rm ref})_{33}-R_{33}\rvert>0.8\}.
\end{equation}
第二项是两种根坐标系中重力方向z分量差的绝对值，不是欧拉角差；参考直立时阈值对应约78.5度倾斜。代码没有末端执行器偏差终止，与论文附录A一致。

上下文 $\xi=[\mu,\Delta q^0,\Delta c_{\rm torso}]\in\R^{33}$ 包含摩擦、29维关节零位偏差与3维躯干质心偏移，分布范围与表A.3相同。摩擦和质心偏移在环境启动时采样一次，关节偏移每个 episode 重采样；所以“每个 episode 的 context”不表示每次 reset 都重采样所有分量。训练 reset 还叠加位置、姿态、速度、关节角扰动，不能和持续施加外部推力混为一谈。

FK 沿关节树递推 $p_i=p_{\rm par}+R_{\rm par}\bar p_i$、$R_i=R_{\rm par}\bar R_i\exp([u_iq_i]_\times)$，并传播 $\omega_i=\omega_{\rm par}+u_i^{\rm world}\dot q_i$、$v_i=v_{\rm par}+\omega_{\rm par}\times(p_i-p_{\rm par})$。这些操作由 PyTorch 构成，允许身体回报对预测状态求导。\textbf{FK是可微运动学，不是可微接触物理引擎。} 接触影响由动力学网络从仿真转移中学习。
\source{\src{supermimic/mdp.py}{147}{161}；\src{supermimic/env.py}{72}{90}；\src{supermimic/env.py}{142}{184}；\src{supermimic/kinematics.py}{62}{90}}

\FloatBarrier
\section{世界模型：从论文式(1)到实际误差函数}\label{sec:wm}
\paper{第3.2节、式(1)（第5页）；表A.4、A.5（第15--16页）}
\subsection{动力学与解析辅助头的职责划分}
论文允许世界模型的奖励、观测、终止头既可以学习，也可以解析计算。仓库选择只学习
\begin{equation}
 \hat S_{t+1}=f_\omega(\hat S_t,a_t,\xi),
\end{equation}
其余由 \code{MDP} 和参考动作库计算：回报依赖预测状态、当前/上一动作与参考帧；观测依赖预测状态、上一动作、零位偏差及参考根朝向；终止依赖预测状态与参考帧。因而没有独立 reward/observation/termination 网络，也没有这些头的 MSE 或 BCE 监督项。\textbf{这不是漏写论文式(1)的损失}，而是论文明确许可的解析头分支。

模型输入不包含参考动作，说明它拟合的是统一机器人动力学，而不是“第几段动作的下一姿态”。$\xi$ 中的摩擦、标定与质心参数使同一确定性网络可以条件化于不同物理系统；传感器噪声由外部采样，不要求网络把噪声本身当成不可预测动力学。

\subsection{局部特征、网络与状态积分}
令 $H=\heading(R)$，输入的 heading/xy 不变特征为
\begin{equation}
 \phi(S)=\big[p_z,\six(H^\top R),R^\top v,R^\top\omega,q,\dot q\big]\in\R^{71}.
\end{equation}
它移除全局xy位置和 yaw，保留高度、倾斜、机体系速度及关节状态。上下文先按固定量纲归一化：$(\mu-0.75)/0.45$、$\Delta q^0/0.01$、$\Delta c/0.05$；再将 $[\phi(S),a,\xi_n]$ 送入在线 \code{RunningNorm}。

动力学网络隐藏宽度为 $(2048,2048,2048)$，每个隐藏层采用 Linear、LayerNorm、SiLU，输出70维标准化局部增量。输出经运行均值和标准差反标准化后拆为
\begin{equation}
 \Delta=(\Delta p_h,\Delta\vartheta_b,\Delta v_b,\Delta\omega_b,\Delta q,\Delta\dot q).
\end{equation}
其积分规则是
\begin{align}
 p'&=p+H\Delta p_h,& R'&=R\exp([\Delta\vartheta_b]_\times),\\
 v'&=v+R\Delta v_b,&\omega'&=\omega+R\Delta\omega_b,\\
 q'&=q+\Delta q,&\dot q'&=\dot q+\Delta\dot q.
\end{align}
位移用 yaw 坐标系增量，速度与旋转增量用根机体系；不能把整个70维输出都称为同一局部坐标系。使用 $SO(3)$ 指数映射避免直接加旋转矩阵元素导致的明显非正交问题。

最后线性层权重和偏置初始化为零。但由于输出还要反标准化，运行均值更新后，零网络输出表示\textbf{平均增量}，并非始终严格的恒等转移。上述特征、LayerNorm、激活函数与初始化是代码实现选择；论文表A.4只规定了主要隐藏层宽度，未逐项给出这些细节。
\source{\src{supermimic/mdp.py}{63}{81}；\src{supermimic/mdp.py}{183}{184}；\src{supermimic/networks.py}{12}{21}；\src{supermimic/networks.py}{67}{113}}

\subsection{多步自回归训练与真实损失}
论文的抽象损失为
\begin{equation}
 \mathcal L_{\rm WM}=\E_{\mathcal B}\sum_{j=1}^{k}\Big[
 \|\hat s_j-s_j\|^2+(\hat r_j-r_j)^2+
 \|\hat{\bar o}_j-\bar o_j\|^2+\operatorname{BCE}(\hat d_j,d_j)\Big],
\end{equation}
解析头对应的项不参与学习。代码的具体操作为
\begin{equation}
 \hat S_0=S_0,\qquad
 \hat S_{j+1}=f_\omega(\hat S_j,a_j^{\rm logged},\xi_j),\quad j=0,\ldots,7.
\end{equation}
注意除了第一帧，\textbf{后续不再用真实 $S_j$ 覆盖预测}。这才是自回归窗口拟合，与把8个独立单步样本打包训练不同。

误差向量也不是71维平铺状态直接相减，而是70维局部几何误差：
\begin{equation}
 e(\hat S,S)=\left[H^\top(\hat p-p),\log(R^\top\hat R)^\vee,
 R^\top(\hat v-v),R^\top(\hat\omega-\omega),\hat q-q,\hat{\dot q}-\dot q\right].
\end{equation}
令 $\sigma_\ell$ 为真实单步增量的运行标准差，$\epsilon_\ell$ 为分量下界，则\emph{实际用于反向传播}的损失为
\begin{equation}\label{eq:actualwm}
 L_{\rm code}=\frac1B\sum_{b=1}^{B}\sum_{j=1}^{8}m_{bj}
 \frac1{70}\sum_{\ell=1}^{70}
 \left[\frac{e_\ell(\hat S_{bj},S_{bj})}{\max(\sigma_\ell,\epsilon_\ell)}\right]^2.
\end{equation}
六类分量下界分别是 $0.003,0.005,0.02,0.05,0.003,0.05$，对应位置、旋转、线速度、角速度、关节角、关节速度。$m_{bj}$ 在窗口越过 episode 边界后变0。

这里有三个容易忽略的实现差别：一是按维度标准化后的误差，而不是无量纲处理前的欧氏和；二是按 batch 数 $B$ 除，不按有效时间步数除，所以短有效窗口的贡献更少；三是日志 \code{wm_loss} 再除窗口长度8，但反向传播用的是式\eqref{eq:actualwm}，不能用日志数值直接重建优化尺度。
此外，日志 \code{wm_err8} 是窗口最后一步在有效样本上的标准化状态误差平均；论文图3(c)则按预定义的 nominal error 归一化。现有材料不足以认定两者使用同一误差标度，不能只看日志数值接近就宣布模型精度复现。
\source{\src{supermimic/train.py}{90}{108}；\src{supermimic/train.py}{147}{156}；\src{supermimic/networks.py}{115}{134}}

\subsection{为什么这能提供梯度，以及不能保证什么}
在 actor 阶段，冻结 $\omega$ 不妨碍计算 $\partial f_\omega/\partial S$ 与 $\partial f_\omega/\partial a$；这些雅可比构成策略 BPTT 的关键链路。8步监督鼓励网络处理自身滚动误差，32步想象则提出更长预测要求。\textbf{8步拟合好并不在数学上保证32步梯度准确，更不能推出无限长稳定性}；论文对长 rollout 稳定性的叙述属于其实验观察，代码与本次 CPU 检查不构成同样证据。

此外，预测误差小不等价于状态/动作雅可比精确，策略仍可能利用模型偏差。短 horizon、critic、持续在线仿真采样和多步拟合有助于降低这些风险，但这里没有模型集成、不确定性惩罚或显式信赖域，不能擅自补充为“论文里隐含实现了这些机制”。

\FloatBarrier
\section{策略与价值学习：论文式(2)、式(3)的梯度实现}\label{sec:actor}
\paper{第3.3节、式(2)--(3)（第5页），附录B--D（第14--16页）}
\subsection{策略网络并不是直接做行为克隆}
策略采用参考编码器加观测解码器：
\begin{align}
 z_t&=\tanh\big(\operatorname{enc}_\theta(N_g(g_t))\big)\in[-1,1]^{32},\\
 a_t&=\operatorname{dec}_\theta\big(z_t,\operatorname{flat}(N_o(o_{t-9:t}))\big).
\end{align}
编码器与解码器隐藏宽度均为 $(2048,1024,512,256)$，默认激活为 ELU；输出层为线性层。解码器最后一层初始权重再乘0.01、偏置置零，以减小初始动作。输入归一化的均值/方差作为 module buffer 保存在 checkpoint 中，而非学习参数。

虽然目标来自参考动作，策略损失里没有 $\|a-a^{\rm ref}\|^2$ 或参考动作的最大似然项。参考定义的是未来目标和跟踪奖励；动作通过模型产生状态、状态再产生回报，策略因此通过控制优化学习，而不是直接复制一个给定的电机动作标签。
\source{\src{supermimic/networks.py}{137}{165}；\src{supermimic/train.py}{76}{81}}

\subsection{想象轨迹的初始化与逐步推进}
\code{sample_valid_starts()} 从 replay 抽起点及其真实历史，并拒绝不够剩余参考帧数的候选，要求 $f+H\le f_{\rm end}$。这避免32步想象越过 clip 尾部。起始状态来自数据、无需梯度；历史则从缓存的状态、上一动作、上下文和帧号重新构造无噪观测，再重新加传感器噪声，而非保存并重用当时的每一次噪声样本。

每个想象步实际顺序为
\begin{enumerate}
\item 用当前历史 $o_{t-9:t}$ 与目标 $g_t$ 计算 $a_t$；
\item 用固定世界模型求 $\hat S_{t+1}$，参考索引推进为 $f+1$；
\item 在新状态上构造观测并加入噪声，将历史向前滚动；
\item 查询新参考索引的未来10帧目标，供下一步动作使用；
\item 最终批量计算 $r_{t+1}=r(\hat S_{t+1},f+1,a_t,a_{t-1})$ 与软终止。
\end{enumerate}
这一时序使 \code{As} 保存的是 $a_t$、\code{Ss} 保存的是 $S_{t+1}$、\code{Aps} 保存的是旧 $a_{t-1}$。如果把奖励错误地对齐到旧状态或旧参考帧，虽然形状可能一致，算法语义已不相同。
\source{\src{supermimic/train.py}{111}{125}；\src{supermimic/train.py}{159}{169}；\src{supermimic/train.py}{175}{200}}

\subsection{软终止、存活权重与有限时域目标}
代码没有学习 termination probability，而是把硬阈值平滑化：
\begin{equation}
 \hat d_{t+1}=1-
 \left[1-\sigma\!\left(\frac{e_z-0.25}{0.02}\right)\right]
 \left[1-\sigma\!\left(\frac{e_{\rm tilt}-0.8}{0.04}\right)\right].
\end{equation}
两个分母是实现选择，原论文未规定这些温度。该构造令接近失败边界的状态已经有生存折损，同时允许存活概率对状态求导。

定义 $\hat c_0=1$，$\hat c_{t+1}=\hat c_t(1-\hat d_{t+1})$，则 actor 最大化
\begin{equation}\label{eq:actor}
 \hat J(\theta)=\E\left[\sum_{t=0}^{H-1}\gamma^t\hat c_t\hat r_{t+1}
 +\gamma^H\hat c_H V_\psi(\hat S_H,g_H,\xi)\right],\quad H=32.
\end{equation}
代码先以 \code{cumprod(1-d,0)} 得到 $c_{t+1}$，再在前面补一行全1并去掉最后一行，得到乘奖励的 $c_t$。因此触发终止的那个转移的即时奖励仍按“进入该步之前的存活权重”计算；不能直接把 $c_{t+1}$ 乘到同一个 $r_{t+1}$，否则与论文式(2)不同。

32步对应0.64秒物理时间。末端 critic 估计更远未来的价值，避免必须完整反传数秒到数分钟的动作。\code{--no_critic} 会去掉该 bootstrap 与后续价值训练，可用于对比有限时域目标，但仓库没有提供论文全部 horizon 消融结果。
\source{\src{supermimic/mdp.py}{154}{161}；\src{supermimic/train.py}{199}{210}}

\subsection{“冻结参数”不等于“切断梯度”}
代码先对 world model 和 critic 参数调用 \code{requires_grad_(False)}，再正常执行其前向；对 actor 损失调用 \code{backward()} 后，只更新策略。它没有把整个想象过程放进 \code{torch.no_grad()}，也没有逐步 \code{detach()} 预测状态。因此保留了
\begin{equation}
 \theta\longrightarrow a_t\longrightarrow\hat S_{t+1}
 \longrightarrow\{\hat r_{t+1},\hat d_{t+1},o_{t+1},a_{t+1},V_H\}
 \longrightarrow\hat J
\end{equation}
这条链路。末端 $V_H$ 虽然参数冻结，其\textbf{输入梯度}仍会通过 $S_H$ 回到 actor；这不是把一个常数 baseline 加到目标中。

用 $x_t$ 表示包含预测状态与可微观测历史的扩展状态，则链式法则具有
\begin{equation}
 \frac{\dd x_{t+1}}{\dd\theta}=
 \frac{\partial F}{\partial x_t}\frac{\dd x_t}{\dd\theta}
 +\frac{\partial F}{\partial a_t}
 \left(\frac{\partial\pi_\theta}{\partial\theta}
 +\frac{\partial\pi_\theta}{\partial x_t}\frac{\dd x_t}{\dd\theta}\right)
\end{equation}
的结构。参考帧索引与抽样噪声不对策略参数求导；给定噪声样本，观测加噪操作仍保留对预测状态的导数。actor 优化的是\textbf{模型中有限样本回报的解析梯度}，不是对真实仿真系统回报的无偏梯度保证。

\begin{caution}{不要用 PPO 的术语误读这段实现}
这里没有随机动作分布的 log-prob、old policy 概率比、clipped surrogate 或 entropy bonus。缓存虽是离策略数据，但 actor 用当前策略在模型内重新 rollout。GAE 在这里用于 critic 的想象回报目标，不意味着 actor 在做 PPO。
\end{caution}

策略最小化 $-\hat J$；反向后先全局范数裁剪至1.0，再逐元素裁剪至0.1。代码先检查裁剪前返回的范数是否有限，只有有限时才执行 actor optimizer step。世界模型和 critic 参数随后恢复可训练状态。
\source{\src{supermimic/train.py}{185}{222}}

\subsection{critic 输入与固定的 \texorpdfstring{$\lambda$}{lambda}-return 目标}
critic 隐藏层是 $(256,128)$，输入包含71维不变状态特征、29维参考关节角误差、14个身体的位置误差42维、根位置误差3维、相对根朝向6维、上下文33维，以及730维目标，总计914维。这些特征在想象状态上解析计算，不是 actor 同款的990维噪声历史。

actor rollout 后，用 \code{no_grad()} 与 \code{detach()} 形成固定的 $H+1$ 个 critic 输入状态，更新 critic 的归一化统计，计算旧 $V_0,\ldots,V_H$。从 $\hat R_H=V_H$ 开始倒序递推
\begin{equation}\label{eq:lambda}
 \hat R_t=\hat r_{t+1}+\gamma(1-\hat d_{t+1})
 \big[(1-\lambda)V_{t+1}+\lambda\hat R_{t+1}\big],\qquad\lambda=0.95.
\end{equation}
固定这组 $\hat R_t$，最小化均方误差 $\E[(V_\psi(x_t)-\hat R_t)^2]$。所有 $H\times B=131,072$ 个时间批样本展平，每轮打乱后分4批，训练5轮，共20次 critic 参数更新。actor 使用的 bootstrap 值来自这轮 critic 更新\emph{之前}。

“固定世界模型”与“固定 critic targets”是两层不同的 stop-gradient：前者冻结动力学权重而允许状态梯度；后者对价值目标及采集的想象特征断开图，避免 critic 拟合再次改变 actor 或世界模型。代码没有单独 target critic 或 Polyak 网络。
\source{\src{supermimic/mdp.py}{163}{180}；\src{supermimic/train.py}{223}{252}}

\FloatBarrier
\section{经验回放与重置采样：训练分布如何形成}
\paper{第3.1节，附录D.1（第4、15页）}
\subsection{环形缓存存了什么，没存什么}
\code{ReplayBuffer} 按 $(T,N)$ 布局直接存 GPU 张量，默认 $T=1000,N=4096$。每个转移包括 \code{s,s1,a,a_prev,xi,r,f,ep_step,done}。存储的动作是\textbf{实际执行的加噪动作}，不是策略加噪前的均值。否则模型输入与仿真产生的转移不对应，监督将出现系统性错配。

字段 \code{done} 是硬失败、10秒超时或 clip 结束的并集。\code{TrackingSim.step()} 虽还返回 \code{term} 与分项回报，但 \code{FIELDS} 不保存它们。由于辅助头解析计算，模型拟合不需要缓存 \code{term} 的 BCE 标签；\code{done} 在这里主要是序列边界。

每个转移的主存储约为
\begin{equation}
 (71+71+29+29+33+1)\times4+8+4+1=949\ \text{bytes},
\end{equation}
因此默认缓存主张量约 $3.89\times10^9$ bytes，即3.62 GiB，不含网络、优化器、临时窗口、想象计算图与仿真器开销。这是由字段形状估算的主张量容量，\textbf{不是实测训练峰值显存}。
\source{\src{supermimic/buffer.py}{9}{39}；\src{supermimic/env.py}{203}{213}}

\subsection{窗口掩码与历史恢复的两种边界处理}
模型窗口先按时间和环境抽取连续转移，再设置
\begin{equation}
 m_j=\prod_{i=0}^{j-1}(1-\mathrm{done}_i),\qquad m_0=1.
\end{equation}
如果窗口第2个转移结束 episode，mask 为 $[1,1,0,\ldots]$；终止转移本身仍可学习，reset 后的另一段 episode 不与前一段拼接拟合。存下来的 $s_1$ 是 reset \emph{之前} 的真实后继状态；reset 后的新状态出现在后续行。

actor 起点需要10帧历史。\code{sample_starts()} 的历史向后距离同时受 episode 已进行步数与缓存可用年龄限制。reset 之前或缓存最早行之前的历史使用最早可用帧重复填充，而不是读取物理不连续的旧 episode。环形存储的物理行号与时间年龄通过 \code{_row()} 转换，避免 wrap-around 后把最新和最旧当作邻接历史。

\code{sample_valid_starts()} 对参考尾部采用拒绝采样；在默认长动作库中有可用起点，但函数没有“找不到足够长 clip”的退出条件。若替换为全都短于想象 horizon 的动作，或可用缓存只有尾帧，应先验证 $H$ 与数据长度。这是代码边界检查的缺口，不是本报告已观察到默认训练卡住。
\source{\src{supermimic/buffer.py}{41}{70}；\src{supermimic/train.py}{159}{169}}

\subsection{基于回报的困难片段采样}
动作时间线被划为1秒、即50控制帧的 bins。设总 bin 数为 $M$，$G_b$ 为该 bin 的回报估计，$G_{\max}=0.1/(1-0.99)=10$。当前困难度为
\begin{equation}
 w_b=\frac{\exp(-25G_b/G_{\max})}{\sum_{b'}\exp(-25G_{b'}/G_{\max})},\quad
 P(b)=\frac{0.5}{M}+0.4w_b+0.1\frac{\bar w_b}{\sum_{b'}\bar w_{b'}}.
\end{equation}
负号意味着低回报 bin 更容易被采样；0.5均匀项避免只盯困难动作，0.1历史项帮助已学会的难例保持练习。代码的历史软权重采用 $\bar w\leftarrow0.999\bar w+0.001w$，具体EMA系数属于实现选择。

重置后从参考状态开始，加随机扰动。回报仅在 episode 结束时分配给\emph{该 episode 起始 bin}，不是给沿途经过的所有 bin。多个环境在同一次结束于同一 bin 时，\code{update_bin_returns()} 先求这些回报的均值，再做
\begin{equation}
 G_b\leftarrow0.8G_b+0.2\,\operatorname{mean}\{G^{(i)}:\mathrm{startbin}_i=b\}.
\end{equation}
这避免同一 bin 的重复索引写入导致与环境顺序有关的更新。
\source{\src{supermimic/motion.py}{72}{87}；\src{supermimic/env.py}{131}{140}；\src{supermimic/reset_sampling.py}{6}{13}}

\subsection{超时 bootstrap 与 critic bootstrap 不是一回事}
训练 episode 最多500步（10秒），或运行到 clip 尾部；硬失败则立即终止。为更新重置采样的 $G_b$，代码对非硬失败结束添加
\begin{equation}
 G=\sum_{t=0}^{T-1}\gamma^t r_{t+1}
 +\mathbf1_{\rm not\ terminal}\,\gamma^T
 \frac{\frac1T\sum_{t=0}^{T-1}r_{t+1}}{1-\gamma}.
\end{equation}
这是用\textbf{本段平均奖励率}近似无限后续的启发式，不是调用 critic 估值；原论文没有给出该具体截断处理。想象 actor 的 bootstrap 则是式\eqref{eq:actor}中的 $V_\psi$。两者服务不同目的，不能用一句“所有超时都用 critic bootstrap”概括。
\source{\src{supermimic/env.py}{207}{235}}

\FloatBarrier
\section{运行工程与超参数：算法之外的必要语境}
\subsection{物理后端和时间步}
代码运行在 Isaac Lab 3 的 kit-less 路径，通过 Newton 的 MuJoCo-Warp solver 做并行仿真，不是默认 Isaac Sim/Kit/PhysX 工作流。\code{tools/setup.sh} 将 Isaac Lab 固定到
\begin{center}\small\code{46d6cc3d7b115906fad12efd372fbf69f4e3d3ed}\end{center}
再构建镜像。控制周期为0.02秒，每步4个物理子步，即5 ms；solver 配置采用 \code{implicitfast} 积分、pyramidal friction cone、MuJoCo contact pipeline，约束与接触容量设置为 \code{njmax=500,nconmax=128}。

G1资产经 URDF $\to$ MJCF $\to$ USD 转换后，用 \code{Physics=mujoco} variant 加载；动作CSV、机器人资产与 Isaac Lab 均不随仓库完整附带。它们的版本和转换结果会影响关节顺序、碰撞体与物理行为，所以“只 clone Python 代码”不是完整可运行环境。
\source{\src{supermimic/env.py}{30}{45}；\src{supermimic/robot.py}{133}{146}；\src{tools/setup.sh}{1}{11}；\src{tools/download_data.sh}{1}{16}}

\Needspace{10\baselineskip}
\subsection{与论文表A.5逐项核对的默认参数}
\begin{longtable}{@{}P{45mm}P{31mm}P{78mm}@{}}
\caption{固定版本 \code{train.py} 的默认配置及其解释。}\label{tab:hyper}\\
\toprule 参数 & 代码默认值 & 对齐情况或注意事项\\\midrule\endfirsthead
\toprule 参数 & 代码默认值 & 对齐情况或注意事项\\\midrule\endhead
并行仿真环境 & 4096 & 对齐论文单GPU配置\\
回放缓存 & 1000步/环境 & 共约409.6万条转移，不是总共1000条\\
模型窗口 / 想象时域 & 8 / 32 & 对齐；分别为0.16 / 0.64秒\\
未来目标 / 观测历史 & 10 / 10帧 & 主体长度对齐；观测和目标的具体量由代码实现\\
模型批次 & $10\times2048$ & 每次外层迭代10次优化，每批2048个窗口\\
actor batch & 4096条想象 & 默认共131,072个想象转移\\
critic 更新 & 5轮 $\times$ 4批 & 共20次optimizer step\\
学习率 & $2\times10^{-4}$ & 三个优化器初值一致\\
优化器 & Adam & actor $\beta=(0.9,0.95)$；其他使用默认beta\\
折扣 / GAE & 0.99 / 0.95 & 对齐表A.5\\
动作采样噪声 & $\mathcal N(0,0.1^2)$ & 对齐；只用于仿真收集，不用于想象动作\\
梯度裁剪 & norm 1.0 & actor 再逐元素裁剪0.1\\
update-to-data ratio & 0.4 & 以外层迭代/向量化仿真步计，不以单环境transition计\\
学习率日程 & cosine $\to0.03\times$ & 250,000次外层迭代后保持常数\\
最大训练迭代 & 320,000 & 默认上限；作者记录的LAFAN1成功点是14,000\\
初始数据收集 & 64个向量步 & 在首次更新前；代码额外选择\\
仅模型预热 & 500次迭代 & 开始阶段不做actor/critic更新；代码额外选择\\
MLP 数值精度 & bf16 autocast & 不同于论文TF32设定；状态与解析头保持fp32\\
多GPU & 无相应训练实现 & 默认固定 \code{cuda:0}，没有DDP/梯度汇总逻辑\\
\bottomrule
\end{longtable}
\source{\src{supermimic/train.py}{18}{45}；\src{supermimic/train.py}{61}{81}；\src{supermimic/train.py}{135}{138}；\src{supermimic/train.py}{290}{321}}

\subsection{UTD 的实际调度与计算量}
代码每轮累加 $1/0.4=2.5$ 个向量化仿真步，借助小数累加器交替执行2步和3步。一次采样步含4096个环境，所以每外层迭代平均新增10,240条物理转移。之后进行10次模型 minibatch 更新、一次 actor 更新和20次 critic 更新。\textbf{这里的0.4不是“模型梯度步数/单条物理转移数”}，也不是把所有21或31次 optimizer step 一起计入的比例。

\code{RunningNorm} 在真实采样数据上持续更新输入/输出及策略观测/目标的统计；critic 的统计从想象特征更新。达到约 $2\times10^8$ 个样本后停止统计更新。世界模型输出尺度来自\emph{真实单步增量}，而非全状态的标准差。断点文件保存了这些 buffer，所以加载模型时应连同归一化状态一起使用。
\source{\src{supermimic/train.py}{275}{283}；\src{supermimic/train.py}{300}{321}；\src{supermimic/networks.py}{24}{51}}

\subsection{加速策略与其可比性边界}
\begin{figure}[H]
\centering
\includegraphics[width=\textwidth]{figures/paper_fig5_efficiency.pdf}
\caption{论文原图5：想象 rollout 的吞吐与解析梯度的样本效率。引用自 Liao 等（2026）图5，第8页\cite{liao_supermimic_2026}，从 arXiv:2610.28480v1 的 PDF 矢量裁切。左图为 H100/TF32 条件下的世界模型与 MuJoCo Warp 吞吐；右图为 SuperMimic 与 PPO-WM 的想象步数对比。两者均是原论文结果，非当前仓库或本次报告的实测。}
\label{fig:paper-efficiency}
\end{figure}

图\ref{fig:paper-efficiency}(a)讨论\textbf{每秒能产生多少转移}，而(b)讨论\textbf{达到相近回报需要多少想象步}：快速模型与有效梯度是两种不同机制。右图的 PPO-WM 已在同类世界模型内采样，但仍用 PPO 更新，因此不能把图中约64倍的步数差直接解释为仿真引擎加速。对应当前代码，\code{wm.step()} 提供快速可微转移，\code{policy_update()} 则通过模型求回报梯度；仓库没有随附该 PPO-WM 对照的完整训练实现和结果，不能据此宣称已经重做此消融。

\code{torch.compile} 编译世界模型 rollout、误差、想象单步、回报/终止与 critic 输入构造；代码还先堆叠所有想象状态，再批量做 FK 与奖励，减少 Python 调度和小批操作开销。MLP 矩阵乘法用 bf16 autocast，输出转换回 fp32；状态积分、FK和奖励继续使用fp32。训练脚本同时允许TF32，但 bf16 路径不能因此被描述为“严格使用论文的TF32”。

README 自报 bf16 在作者机器上快约24\%，这不是本次跨硬件测量结果。依赖的 GPU 架构、CUDA、PyTorch编译和 solver 都会改变吞吐；没有同环境基准，不应把差异全部归因于算法。
\source{\src{supermimic/networks.py}{54}{64}；\src{supermimic/train.py}{128}{132}；\src{supermimic/train.py}{190}{199}}

\subsection{checkpoint 与 resume 的真实边界}
\code{last.pt} 保存 world model、actor、critic、三个优化器和计数器；编号快照只保存 actor、world model 等子集，用于评估。\code{--resume} 直接读取完整键集合，因此应使用 \code{last.pt}，不能假定评估用的小快照也可无差别续训。

恢复时不载入 replay buffer、物理环境状态、reset bin 回报统计、观测历史或随机数生成器状态；这些在新进程重新创建。所以它是“加载学习状态继续训练”，不是逐位相同的无缝续跑。读取恢复迭代后，\code{lr_at(it)} 会按计数重新设置学习率。

计时起点在环境初始化与最初数据收集之后；构建、下载、资产转换和评估不计入训练日志。500次仅模型优化则在主循环计时内。复现实验应同时报告端到端耗时与该训练计时，避免把两者混为同一种成本。
\source{\src{supermimic/train.py}{82}{87}；\src{supermimic/train.py}{290}{302}；\src{supermimic/train.py}{341}{352}}

\clearpage
\section{论文--代码逐项对齐矩阵}\label{sec:matrix}
下面的“对齐”指算法结构或相应数值配置能够在源码中定位，不表示本次运行证明了相同实验成绩；“实现选择”表示论文允许但没有规定具体形式；“偏离”表示与论文列出的接口或运行条件不同；“未覆盖”表示在该固定版本中未找到相应完整实现。

\small
\begin{longtable}{@{}P{30mm}P{45mm}P{79mm}@{}}
\caption{方法、接口及实验范围的对齐清单。所有源码定位均绑定固定提交。}\\
\toprule 论文对象 & 代码锚点 & 判断及理由\\\midrule\endfirsthead
\toprule 论文对象 & 代码锚点 & 判断及理由\\\midrule\endhead
图2、算法A.1 & \src{supermimic/train.py}{300}{321} & \textbf{主体对齐。} 同一在线循环里收集、拟合、想象学习；开始阶段额外加入模型预热。\\
第3.1节，实际执行动作入缓存 & \src{supermimic/train.py}{274}{283} & \textbf{对齐。} 加噪并裁剪后的动作输入仿真并存入buffer，而非只存确定性动作。\\
第3.2节，条件化动力学 & \src{supermimic/networks.py}{67}{113} & \textbf{主体对齐。} 输入状态特征、动作与33维context；局部特征及增量积分是具体实现选择。\\
式(1)，8步窗口 & \src{supermimic/train.py}{95}{108} & \textbf{结构对齐、度量具体化。} 自回归，mask阻止跨episode监督；误差按自适应尺度标准化。\\
第3.2节，辅助头 & \src{supermimic/mdp.py}{94}{161} & \textbf{允许的实现选择。} 回报、观测、终止解析计算，无额外预测损失；新增参考朝向使观测还依赖目标信息。\\
式(2)，BPTT目标 & \src{supermimic/train.py}{175}{222} & \textbf{主体对齐。} $H=32$、折扣存活奖励、终端价值；只冻结模型和critic权重，保留输入梯度。\\
式(3)，价值目标 & \src{supermimic/train.py}{223}{252} & \textbf{对齐。} 想象数据上倒序$\lambda$-return，目标只计算一次，监督训练5轮$\times$4批。\\
附录A，表A.1观测 & \src{supermimic/mdp.py}{18}{35}；\src{supermimic/mdp.py}{95}{105} & \textbf{偏离。} 93维本体项之外加入6维相对参考朝向和相应噪声，总计99维。不能称完全相同观测的严格受控复现。\\
附录A，未来参考目标 & \src{supermimic/motion.py}{94}{117} & \textbf{长度对齐，内容具体化。} 10个未来帧、每帧73维、730维展平；论文没有完整列出这里的全部特征及尾帧填充。\\
表A.2，回报 & \src{supermimic/mdp.py}{108}{145} & \textbf{主体对齐。} 六跟踪项与权重、动作平滑、关节限位权重一致；坐标重锚定和90\%软限位按源码解释。\\
删除自碰撞奖励 & \src{supermimic/robot.py}{133}{146} & \textbf{需区分两层。} 删除奖励项符合论文；关闭物理自碰撞是仓库额外偏离，不由论文这句话推出。\\
附录A，终止 & \src{supermimic/mdp.py}{147}{161} & \textbf{硬阈值对齐、软头具体化。} 高度0.25m、重力z差0.8；想象sigmoid温度0.02/0.04未由论文规定。\\
表A.3，context/DR & \src{supermimic/env.py}{72}{90}；\src{supermimic/env.py}{182}{184} & \textbf{范围对齐、采样时机不同粒度。} 摩擦/质心每环境启动一次，关节零位每episode。\\
表A.4，网络 & \src{supermimic/networks.py}{67}{77}；\src{supermimic/networks.py}{137}{165} & \textbf{隐藏层宽度及latent对齐。} LayerNorm/激活/初始化属实现选择；实际输入维数随额外特征改变。\\
附录D.1，重置分布 & \src{supermimic/env.py}{131}{140}；\src{supermimic/reset_sampling.py}{6}{13} & \textbf{混合比例对齐。} 50/40/10与softmax系数25一致；EMA、重复bin聚合与截断bootstrap为具体实现。\\
表A.5，训练计划 & \src{supermimic/train.py}{18}{45} & \textbf{大部分对齐。} 4096环境、UTD、batch、学习率与裁剪；额外64步收集及500轮仅模型预热。\\
论文数值精度 & \src{supermimic/train.py}{61}{65}；\src{supermimic/networks.py}{56}{64} & \textbf{偏离。} MLP默认bf16，而论文报告TF32；硬件也不同。\\
附录B/G，多目标接口 & \src{supermimic/networks.py}{137}{165} & \textbf{未覆盖。} 只有motion encoder；没有velocity/末端/SMPL接口、decoder定向冻结和latent-alignment训练。\\
双GPU、规模扩展 & \src{supermimic/train.py}{65}{74} & \textbf{未覆盖。} 固定单设备，无DDP；未提供BONES-SEED过滤/扩展实验全流程。\\
附录F.1，评估 & \src{supermimic/evaluate.py}{52}{96} & \textbf{主要协议对齐。} 首帧确定性、无噪声/DR、相同硬失败；具体回报汇总应区分clip与frame权重。\\
第4.4节，实机转移 & 仓库无对应完整模块 & \textbf{未覆盖。} 仿真评估和视频渲染不构成机器人通信、实时控制与实机验证。\\
\bottomrule
\end{longtable}
\normalsize

论文第4.5节实际展示的下游实验是速度接口的初步研究；附录B提到的末端位姿、SMPL目标属于多接口的泛化描述，不能据此声称原论文逐一完成了这些接口的完整实验。“未覆盖”也不表示这些模块是读懂核心motion-tracking实现的前提。

\begin{caution}{如何解读“非官方复现成功”}
可以说该实现保留了论文核心方法，并由作者记录在其 LAFAN1 设置下达到预设指标。不能据此说它已严格控制所有变量复现论文图4，也不能说其结果证明了论文的全数据规模扩展、多接口控制或实机结论。额外观测、碰撞物理、数值精度及硬件差异都会影响实验可比性。
\end{caution}

\FloatBarrier
\section{评估与性能结论：比较前必须统一分母}\label{sec:evaluation}
\paper{第4.1节、图3--4（第6--7页），附录F.1--F.2（第17--18页）}
\subsection{评估脚本具体做了什么}
\code{evaluate.py} 每个 clip 分配一个环境，全部从该 clip 首帧参考状态启动。评估关闭初始扰动、关节零位与质心随机化，摩擦固定1.0，不加观测噪声和动作噪声。使用 checkpoint 的策略及其归一化统计，推理目标仍是10帧未来参考；\code{pi.eval()} 不等于 \code{no_grad()}，脚本另用装饰器显式关闭梯度。

环境达到硬失败或 clip 结束后，将该环境标为完成；脚本为批处理仍让其物理前进，但累计回报和失败统计只在未完成 mask 下更新，并限制完成环境的帧索引避免越界。这不是失败后 reset 再继续补足同一个 clip 的成绩。

设第 $i$ 个 clip 重采样后有 $L_i$ 帧，回报只累计首次失败/结束前的 $T_i$ 个转移。归一化为
\begin{equation}
 \bar G_i=\frac{\sum_{t=0}^{T_i-1}r_{i,t+1}}{(L_i-1)\,r_{\max}},\qquad r_{\max}=0.1.
\end{equation}
分母始终是\textbf{原 clip 完整转移数}，而不是活过的 $T_i$。这样提前失败后没有继续获得的奖励也反映在分数里。该回报无时间折扣，且包含负正则项，因此不是纯姿态误差或成功率，也不保证下界为0。
\source{\src{supermimic/evaluate.py}{52}{96}；\src{supermimic/env.py}{64}{66}；\src{supermimic/env.py}{215}{223}}

脚本同时输出两种聚合：
\begin{align}
 \texttt{norm\_return}&=\frac1C\sum_{i=1}^{C}\bar G_i,\\
 \texttt{norm\_return\_framew}&=
 \frac{\sum_i\sum_t r_{i,t+1}}{\sum_i(L_i-1)r_{\max}}.
\end{align}
前者每个 clip 等权，后者每个参考转移等权，长 clip 权重更高。论文附录F.1明确了首帧、摩擦、无噪声与失败判据，但没有在该处完全展开这两种聚合的选择。仓库自报数字应按它实际写入的 \code{norm_return} 理解；不能在不同汇总分母下直接宣布严格复现同一曲线。

\subsection{论文原图：规模扩展与 LAFAN1 对照}
\begin{figure}[H]
\centering
\includegraphics[width=\textwidth]{figures/paper_fig3_scaling.pdf}
\caption{论文原图3：跟踪回报、训练成本与数据规模。引用自 Liao 等（2026）图3，第7页\cite{liao_supermimic_2026}，从 arXiv:2610.28480v1 的 PDF 矢量裁切。原图(a)比较 SONIC 与 SuperMimic 的 GPU-hours；(b,c)展示1至264小时动作库的成本与8步预测误差。保留原坐标轴、图例与预算标记；不是40段 LAFAN1 仓库的训练曲线。}
\label{fig:paper-scaling}
\end{figure}

图\ref{fig:paper-scaling}(a)的主要基准是129,514段 BONES-SEED 动作；(b,c)刻画同一机器人动力学在不同数据覆盖下的扩展行为。它支持“一个模型可共享于很多动作”的论文动机，但当前代码没有提供这一数据扩展实验的完整复现链路。特别是(c)的 nominal-error 归一化与代码 \code{wm_err8} 的自适应标准化不能直接视为同一标尺，硬件与训练数据规模的差异也应一并保留。

\begin{figure}[H]
\centering
\includegraphics[width=\textwidth]{figures/paper_fig4_controlled.pdf}
\caption{论文原图4：LAFAN1 受控比较与 critic/horizon 消融。引用自 Liao 等（2026）图4，第7页\cite{liao_supermimic_2026}，从 arXiv:2610.28480v1 的 PDF 矢量裁切。所有子图使用40段、约2.45小时的 LAFAN1 动作；(a)是 checkpoint 在完整动作库上的评估回报，(b,c)使用训练回报。阴影标出默认 $H=32$，叉号表示未达目标。}
\label{fig:paper-controlled}
\end{figure}

图\ref{fig:paper-controlled}(a)才是该仓库 README 声明的直接对照目标；(b,c)则解释了默认 \code{--horizon 32} 及末端 critic 的设计背景。代码提供 \code{--no_critic} 和可配置 horizon，但这不等于已交付完整消融结果。还需注意：训练 episode 采用基于回报的重置分布，而(a)从每个 clip 首帧评估，所以\textbf{不能把(b,c)的训练回报数值与(a)的评估分数逐点等同}。图中成绩均归原论文，仓库作者记录另在下表单独列出。

\subsection{三个层次的成绩不应混在一张结论里}
\begin{table}[H]
\centering\small
\noindent\begin{tabularx}{\textwidth}{@{}P{30mm}P{40mm}Y@{}}
\toprule
来源/设置 & 报告的训练成本 & 指标与适用范围\\
\midrule
论文LAFAN1，图4(a) & 1张H100，1.2h & 40段、约2.45h；归一化回报达到0.82\\
仓库LAFAN1，里程碑 & 1张RTX PRO 6000，1.24h & 迭代12,000；回报0.824；仍有1/40失败\\
仓库LAFAN1，里程碑 & 同GPU，1.47h & 迭代14,000；回报0.829；0/40失败\\
论文BONES-SEED & 1张H100，3h & 129,514段、264.56h；25段失败\\
论文BONES-SEED & 2张H100，9.8h & 同主基准；0段失败；成本19.6 GPU-h\\
本次报告 & 未进行训练 & 16项CPU检查通过；不产出策略回报或失败数\\
\bottomrule
\end{tabularx}
\caption{论文成绩与非官方仓库自报成绩分开列示。仓库数字来自 README 与 MILESTONES，不是本次测量。}
\end{table}
\source{\src{README.md}{1}{24}；\src{MILESTONES.md}{1}{23}}

论文将 SuperMimic 的3 GPU-h与 SONIC 的 $128\times168=21,504$ GPU-h比较，数量级比约为7,168；但论文也说明GPU型号、训练数据规模和可用峰值算力不同。该计算是理解论文预算口径，不是这个40段仓库已重新测得“千倍加速”。同理，作者1.47小时是单GPU训练时间记录，不包含全部构建、下载和评估成本，也不能当作任意硬件的必达承诺。

\subsection{成功率、跟踪质量与泛化是三种问题}
失败判据只有高度和重力z分量差。\textbf{0/40失败不表示每一关节都准确匹配}，也不表示没有任何物理异常；还需看回报、分项误差、动作平滑和具体动作轨迹。

论文附录F明确，其主基准是训练动作上的覆盖率，而非未见动作泛化。该仓库也在训练所用LAFAN1动作库上评估，不能把40/40成功写成“对未见动作100\%泛化”。速度接口任务和实机迁移属于另一层证据，当前仓库没有提供对应完整复现实验。

\subsection{当前公开材料足以支持哪些判断}
固定提交包括结果曲线图、演示GIF和文字里程碑，但没有随版本跟踪完整训练CSV与checkpoint。因此可以审查实现是否具有论文所述机制，亦可准确转述作者的结果记录；不能只凭图像和README独立审计所有训练时长、每段动作指标或训练随机性。

若进一步做严格实验，应固定数据文件校验和与资产转换版本，记录实际命令、GPU、驱动、CUDA 与 PyTorch、碰撞开关和精度，保留每clip回报/失败信息，并分别报告clip-mean与frame-weighted分数。若要声称观测严格对齐，应先单独消融新增6维参考朝向，而不是悄悄删除它后仍沿用已有成绩。

\FloatBarrier
\section{本次验证、局限与接手建议}\label{sec:verify}
\subsection{真正运行过的检查}
随附 \code{verify_alignment.py} 要求目标源码的 HEAD 与固定提交相同，对所有上游 Python 文件进行 AST 语法解析，并运行下表测试。环境为 Python 3.11.18、PyTorch 2.15.0+cpu、NumPy 2.4.6，结果\textbf{16/16通过}。

\noindent\begin{tabularx}{\textwidth}{@{}P{30mm}P{10mm}Y@{}}
\toprule
类别 & 数量 & 实际检查内容\\
\midrule
上游buffer测试 & 2 & 空缓存错误信息；恰好8步缓存可抽8步窗口\\
上游reset测试 & 3 & 重复bin先平均再EMA；顺序无关；空批不变\\
上游goal测试 & 3 & 正常参考特征；跨尾帧保持姿态且仅补齐部分速度置零；不修改原数据\\
新增接口检查 & 2 & 状态编解码及维数；heading/xy不变特征在变换前后一致\\
新增动力学检查 & 2 & 小宽度网络受控输出下的局部增量积分；冻结模型权重仍有动作梯度\\
新增回放检查 & 2 & 终止转移保留但后续被mask；历史不跨episode/环形缓存起点\\
新增目标检查 & 2 & 存活权重索引与独立$\lambda$递推计算；软终止值与梯度方向\\
\bottomrule
\end{tabularx}

为绕开未安装 Isaac Lab 的导入链，脚本显式替换了 \code{supermimic.kinematics} 的\textbf{导入入口}，仅供不调用FK的CPU张量逻辑使用。它不提供伪造的物理结果，不测试FK正确性。动力学积分检查使用原始 \code{WorldModel.step()}、较小隐藏层与人为设定输出；证明对应积分路径可工作，不证明训练后的大模型预测准确。存活/GAE检查是独立参考计算，未调用依赖 Isaac Lab 的完整训练主循环。

\begin{caution}{验证范围不能越界}
本次未运行 Newton/MuJoCo-Warp、资产转换、真实FK比对、CUDA、bf16、\code{torch.compile}、训练收敛、checkpoint策略评估、多GPU或实机控制。静态定位和CPU测试均不能证明论文的性能、物理稳定性与迁移结论。上游README关于FK误差达 $10^{-6}$ 的叙述也未由本次重新测量。
\end{caution}

\subsection{从读懂到严格复现的分层路线}
\begin{enumerate}
\item \textbf{接口层：} 核对29关节顺序、四元数顺序、50 Hz重采样、71/70/99/730/914维契约；运行随附CPU脚本。此层成本低，先排除形状相同但语义不同的问题。
\item \textbf{物理层：} 在仓库固定的Isaac Lab版本内检查资产、FK、接触和动作设定，先少量环境执行现有 \code{test_spawn/test_env/test_reset} 等脚本；不要一上来做长训练。
\item \textbf{学习层：} 分别记录一步/8步/32步预测误差、actor梯度范数、想象存活率与仿真失败率。想象奖励升高而物理回报下降时，优先怀疑模型偏差或接口错配，不能只看 \code{pi_loss}。
\item \textbf{对照层：} 单独测额外6维观测、碰撞开关、bf16/TF32、warmup与误差尺度的影响；每次只改变一个实验变量，保留实际命令和完整指标。
\item \textbf{范围层：} 若目标扩展到论文主实验，再补BONES-SEED过滤与数据组织、多GPU同步、下游接口与相应评估；这些不是改一个 \code{motion_dir} 就自动齐备的功能。
\end{enumerate}

\subsection{最终判断}
该仓库适合作为学习 SuperMimic 核心机制的紧凑入口：三阶段训练循环和三条主要目标函数均能找到具体实现，状态几何、解析回报和BPTT图也便于逐步追踪。最值得掌握的主线是
\begin{center}
\textbf{条件化动力学拟合 $\rightarrow$ 可微想象 $\rightarrow$ 短时域回报梯度 $+$ 价值延伸。}
\end{center}
但“核心方法实现清楚”“在某LAFAN1设置下作者记录达标”与“完整严格复现论文”是不同判断。额外观测、关闭自碰撞物理、数值精度、损失尺度与实验范围都应保留在报告结论中，不能为追求表面一致而省略。

\clearpage
\appendix
\section{源码导航与阅读抓手}
\subsection{按问题而不是按文件长度阅读}
\begin{longtable}{@{}P{41mm}P{77mm}P{36mm}@{}}
\toprule
想回答的问题 & 推荐起点 & 对应论文\\\midrule\endfirsthead
\toprule 想回答的问题 & 推荐起点 & 对应论文\\\midrule\endhead
每轮到底更新什么？ & \src{supermimic/train.py}{290}{321} & 图2、算法A.1\\
真实动作如何进入模型数据？ & \src{supermimic/train.py}{266}{283}；\src{supermimic/env.py}{192}{213} & 第3.1节\\
8步窗口怎样形成？ & \src{supermimic/buffer.py}{41}{55}；\src{supermimic/train.py}{95}{108} & 式(1)\\
局部增量如何变回世界状态？ & \src{supermimic/networks.py}{82}{134} & 第3.2节的实现细化\\
为什么冻结模型还能训练策略？ & \src{supermimic/train.py}{185}{222} & 式(2)\\
GAE目标是否每轮重算？ & \src{supermimic/train.py}{223}{252} & 式(3)，目标每批仅算一次\\
策略是否使用特权参数？ & \src{supermimic/mdp.py}{95}{105}；\src{supermimic/networks.py}{151}{155} & 附录A--B\\
身体奖励如何保持可微？ & \src{supermimic/kinematics.py}{62}{90}；\src{supermimic/mdp.py}{108}{145} & 表A.2\\
未来目标在尾帧如何处理？ & \src{supermimic/motion.py}{94}{120} & Goal的实现选择\\
哪个分布决定困难动作重试？ & \src{supermimic/env.py}{131}{140}；\src{supermimic/env.py}{225}{235} & 附录D.1\\
评估成功和回报怎样定义？ & \src{supermimic/evaluate.py}{69}{91} & 附录F.1\\
为什么checkpoint不等于可精确续训？ & \src{supermimic/train.py}{82}{87}；\src{supermimic/train.py}{341}{352} & 工程边界\\
\bottomrule
\end{longtable}

\clearpage
\subsection{精简伪代码：三个优化器的梯度边界}
以下为解释性伪代码，不是仓库可直接运行的替代脚本：
\begin{lstlisting}
for iteration in schedule:
    # No simulator differentiation; store the action actually executed.
    replay += simulator(policy(noisy_obs, goal) + action_noise)

    for _ in range(10):
        window = replay.sample_windows(B=2048, k=8)
        predicted = autoregress(wm, window.s0, window.logged_actions)
        loss_wm = masked_standardized_state_error(predicted, window)
        step(opt_wm, loss_wm)

    if iteration < 500:
        continue
    start, history = replay.sample_valid_imagination_starts(B=4096)
    freeze_parameters(wm, critic)  # NOT no_grad; retain input Jacobians
    imagined = rollout(wm, policy, start, history, H=32)
    J = discounted_survival_reward(imagined)
    J += gamma**32 * imagined.survival * critic(imagined.last)
    step(opt_policy, -J.mean())
    unfreeze_parameters(wm, critic)

    with no_grad():
        X = detached_critic_features(imagined)
        targets = lambda_returns(X, imagined.r, imagined.d, lam=0.95)
    for epoch in range(5):
        for minibatch in shuffle_into_four(X, targets):
            step(opt_critic, mse(critic(minibatch.X), minibatch.targets))
\end{lstlisting}

\section{复用报告源文件与验证记录}
LaTeX 工程由 \code{main.tex}、\code{sections/}、\code{references.bib} 构成。使用支持中文的 XeLaTeX/BibTeX，或 Tectonic 编译；工程采用CTeX自带Fandol字库，未依赖机器私有字体，引用的论文原图来自同版本 arXiv PDF，按中文图注逐一归属原作者，不把它们当作本次生成的实验结果。

\code{EVIDENCE.md} 记录来源、已执行验证与未验证范围；\code{source_manifest.json} 固定上游文件校验和及PDF校验和；\code{verify_alignment.py} 在外部固定提交的源码目录上运行。\code{figure_sources.json} 和 \code{FIGURES.md} 记录原图编号、页码、裁切坐标、校验和及引用方式；\code{fetch_paper_figures.py} 可重建四张矢量图片。编译产生的报告PDF、图片资产和日志不提交到Git，交付的源码ZIP则包含编译所需图片。

本报告引用少量源码标识和结构，代码来源为MIT许可的非官方仓库，仓库署名 Chong Zhang。数据与机器人资产有独立许可，代码MIT许可不等于数据可自由再分发。本文未向 Zotero 导入新条目、修改原条目或分发动作数据、模型checkpoint与原论文PDF。

\clearpage
\bibliographystyle{unsrtnat}
\bibliography{references}
\end{document}
